\documentclass{article}
\usepackage[main,preprint]{neurips_2026}
\usepackage[utf8]{inputenc}
\usepackage[T1]{fontenc}
\renewcommand{\ttdefault}{pcr}
\usepackage{amsmath,amssymb,graphicx,booktabs,xcolor}
\usepackage[hidelinks]{hyperref}
\usepackage{url,microtype,longtable}
\title{Precision Remaining in a Frozen Score:\\Conditional Gains Beyond Scalar Prediction-Powered Mean Inference}
\author{Anonymous authors}
\date{}
\input{manuscript/numbers}
\begin{document}
\maketitle
\begin{abstract}
Can confidence intervals shrink further when labeled data, unlabeled data and a predictor are fixed? We examine population means using score-only control variates. A variance decomposition explains when nonlinear score information remains after scalar tuning and why affine or uninformative scores offer no first-order gain in this class. Across 165,000 paired draw sets, a cross-fitted cubic reduces mean width by \QuadReduction{} relative to PPI++ in a designed nonlinear Gaussian benchmark, with coverage compatible with 95\%; the three archived-data benchmarks show no practical gain. Rare-event normal intervals severely undercover, and small-label fits can be unstable. Separately, an independently trained bounded-data construction has a finite-sample coverage guarantee and shortens empirical-Bernstein intervals by \BoundedReduction{} versus label-only in a high-signal synthetic setting. It loses to the finite-sample baselines on the archived datasets. The answer is conditional on score information, fitting cost, unlabeled sample size and the required validity regime; no universal interval-length optimum is claimed.
\end{abstract}
\input{manuscript/body}
\clearpage
\bibliographystyle{plainnat}
\bibliography{manuscript/references}
\clearpage
\input{manuscript/appendix}
\input{manuscript/checklist}
\end{document}
